% ===============================================================
% GRADE 8 -- MATATAG Science 8, Three-Term BOW (rev. 8 Apr 2026)
% ===============================================================
\section{Grade 8}
\resetcards
\begin{gradebody}

% ---------------------------------------------------------------
\subsection{Term 1: Body Systems, Heredity, Taxonomy, Energy in Life, Atoms}

\topic{Organ Systems Working Together}
\kf
\begin{keyfacts}
\item Food path: mouth $\rightarrow$ pharynx $\rightarrow$ esophagus $\rightarrow$ stomach $\rightarrow$ small intestine $\rightarrow$ large intestine $\rightarrow$ rectum $\rightarrow$ anus.
\item Processes: ingestion, mechanical processing (chewing, churning), secretion, chemical digestion (enzymes), absorption, elimination.
\item \textbf{Peristalsis}: wave-like muscle contractions that push food along.
\item Saliva has amylase (digests starch). The stomach adds HCl and pepsin (protein). The liver makes bile, which the gallbladder stores; bile emulsifies fats.
\item The small intestine (duodenum, jejunum, ileum) absorbs nutrients through its \textbf{villi}. The large intestine absorbs water.
\item Excretion: kidneys (nephrons) filter blood into urine. The lungs remove \ce{CO2}; the skin removes sweat.
\item Plant transport: roots absorb water, \textbf{xylem} carries water up, \textbf{phloem} carries sugars, and transpiration at the stomata pulls water up.
\end{keyfacts}
\cards
\qa{Wave-like contractions that move food}{Peristalsis}{E}
\qa{Where most digestion and absorption happen}{Small intestine}{E}
\qa{Finger-like projections that absorb nutrients}{Villi}{A}
\qa{Starch-digesting enzyme in saliva}{Salivary amylase}{A}
\qa{Organ that produces bile}{Liver}{E}
\qa{Organ that stores bile}{Gallbladder}{A}
\qa{Role of bile}{Emulsifies fats}{D}
\qa{Protein-digesting enzyme of the stomach}{Pepsin}{A}
\qa{Main job of the large intestine}{Absorbs water}{A}
\qa{Tube from the pharynx to the stomach}{Esophagus}{E}
\qa{Flap that keeps food out of the windpipe}{Epiglottis}{A}
\qa{First part of the small intestine}{Duodenum}{D}
\qa{Functional unit of the kidney}{Nephron}{D}
\qa{Plant tissue that carries water upward}{Xylem}{E}
\qa{Plant tissue that carries sugars}{Phloem}{E}
\qa{Water loss through leaf stomata}{Transpiration}{A}
\begin{confbox}
\confusion{Mechanical}{chemical digestion}{mechanical breaks food into smaller pieces; chemical uses enzymes to break molecules apart.}
\confusion{Digestion}{absorption}{digestion breaks food down; absorption moves nutrients into the blood.}
\confusion{Excretion}{elimination}{excretion removes metabolic waste (urine, \ce{CO2}); elimination (egestion) removes undigested food (feces).}
\confusion{Xylem}{phloem}{xylem carries water up only; phloem carries food both ways.}
\end{confbox}
\mnemonic{\textbf{Ph}loem carries \textbf{f}ood (same ``f'' sound); \textbf{X}ylem carries water up to the sky.}

\topic{Heredity: Patterns of Inheritance}
\kf
\begin{keyfacts}
\item \textbf{Gregor Mendel}, the ``father of genetics,'' worked with garden peas.
\item Alleles are versions of a gene. The dominant allele is written in capitals (T); the recessive in lowercase (t).
\item \textbf{Genotype} is the allele letters (TT, Tt, tt); \textbf{phenotype} is the visible trait.
\item Homozygous: TT or tt. Heterozygous: Tt.
\item Tt $\times$ Tt gives genotypes 1 TT : 2 Tt : 1 tt and phenotypes 3 dominant : 1 recessive.
\item Tt $\times$ tt gives 1 : 1.
\item Laws: Dominance, Segregation, Independent Assortment. A Punnett square predicts the offspring.
\item Pedigree symbols: square = male, circle = female, shaded = affected.
\end{keyfacts}
\cards
\qa{Father of genetics}{Gregor Mendel}{E}
\qa{Mendel's experimental plant}{Garden pea}{E}
\qa{Observable trait}{Phenotype}{E}
\qa{Allele combination (letters)}{Genotype}{E}
\qa{Tt is described as}{Heterozygous}{E}
\qa{Alternative forms of a gene}{Alleles}{A}
\qa{Grid used to predict offspring}{Punnett square}{E}
\qa{Phenotype ratio of Tt $\times$ Tt}{3:1}{A}
\qa{Genotype ratio of Tt $\times$ Tt}{1:2:1}{A}
\qa{Family chart of a trait across generations}{Pedigree}{A}
\qa{Shaded square in a pedigree}{Affected male}{A}
\qa{Tt $\times$ tt: chance of a short (tt) offspring}{50\%}{D}
\qa{tt $\times$ tt: chance of a tall offspring}{0\%}{D}
\qa{Law: allele pairs separate when gametes form}{Law of Segregation}{D}
\begin{confbox}
\confusion{Genotype}{phenotype}{genotype is the letters (Tt); phenotype is what you see (tall).}
\confusion{Homozygous}{heterozygous}{homozygous alleles are the same (TT or tt); heterozygous alleles differ (Tt).}
\confusion{Dominant}{common}{dominant means it shows in Tt. It does \emph{not} mean the trait is most frequent.}
\end{confbox}

\topic{Taxonomic Classification}
\kf
\begin{keyfacts}
\item \textbf{Carolus Linnaeus} is the father of taxonomy. He introduced binomial nomenclature, written \emph{Genus species} (italicized, genus capitalized).
\item Ranks: Domain, Kingdom, Phylum, Class, Order, Family, Genus, Species.
\item Three domains (\textbf{Carl Woese}): Bacteria, Archaea, Eukarya.
\item Six kingdoms: Eubacteria, Archaebacteria, Protista, Fungi, Plantae, Animalia.
\item Humans: Eukarya, Animalia, Chordata, \textbf{Mammalia}, \textbf{Primates}, Hominidae, \emph{Homo sapiens}.
\item Mammals have hair and mammary glands and are warm-blooded. Primates have opposable thumbs, forward-facing eyes, nails, and large brains.
\end{keyfacts}
\cards
\qa{Father of taxonomy}{Carolus Linnaeus}{E}
\qa{Two-name naming system}{Binomial nomenclature}{E}
\qa{Broadest taxon}{Domain}{E}
\qa{Most specific taxon}{Species}{E}
\qa{Taxon between order and genus}{Family}{A}
\qa{Proposed the three-domain system}{Carl Woese}{D}
\qa{The three domains}{Bacteria, Archaea, Eukarya}{A}
\qa{Kingdom of mushrooms and molds}{Fungi}{E}
\qa{Kingdom of amoebas and most algae}{Protista}{A}
\qa{Kingdom of prokaryotes living in extreme conditions}{Archaebacteria}{D}
\qa{Class of humans}{Mammalia}{E}
\qa{Order of humans}{Primates}{E}
\qa{Scientific name of humans}{\emph{Homo sapiens}}{E}
\qa{Two features unique to mammals}{Hair and mammary glands}{A}
\qa{Primate trait for grasping}{Opposable thumb}{A}
\qa{Scientific name of the Philippine eagle}{\emph{Pithecophaga jefferyi}}{D}
\begin{confbox}
\confusion{Domain}{kingdom}{domain is the rank \emph{above} kingdom (3 domains, 6 kingdoms).}
\confusion{Archaea}{Bacteria}{both are prokaryotes, but archaea differ chemically and often live in extreme environments.}
\end{confbox}
\mnemonic{\textbf{D}ear \textbf{K}ing \textbf{P}hilip \textbf{C}ame \textbf{O}ver \textbf{F}or \textbf{G}ood \textbf{S}oup.}

\topic{Photosynthesis, Respiration, and Cycles in Nature}
\kf
\begin{keyfacts}
\item Photosynthesis (in the chloroplast): \ce{6CO2 + 6H2O ->[{light}][{chlorophyll}] C6H12O6 + 6O2}.
\item Cellular respiration (in the mitochondrion): \ce{C6H12O6 + 6O2 -> 6CO2 + 6H2O} + energy (ATP).
\item Stomata exchange gases, and guard cells open and close them.
\item Starch test: iodine turns blue-black. Destarch the plant in the dark before testing.
\item Anaerobic respiration: yeast makes ethanol and \ce{CO2}; muscles make lactic acid.
\item Carbon cycle: photosynthesis removes \ce{CO2}; respiration, combustion, and decay return it.
\item Water cycle: evaporation, transpiration, condensation, precipitation, runoff/collection.
\end{keyfacts}
\cards
\qa{Gas released by photosynthesis}{Oxygen}{E}
\qa{Raw materials of photosynthesis}{\ce{CO2}, water, light}{E}
\qa{Green pigment that absorbs light}{Chlorophyll}{E}
\qa{Leaf pores for gas exchange}{Stomata}{E}
\qa{Cells that open and close the stomata}{Guard cells}{A}
\qa{Formula of glucose}{\ce{C6H12O6}}{A}
\qa{Energy molecule made in respiration}{ATP}{A}
\qa{Reagent to test a leaf for starch}{Iodine solution}{A}
\qa{Colour of iodine with starch}{Blue-black}{A}
\qa{Why a plant is kept in the dark before a starch test}{To use up stored starch}{D}
\qa{Products of yeast fermentation}{Ethanol and \ce{CO2}}{D}
\qa{Muscles short of \ce{O2} produce}{Lactic acid}{D}
\qa{Light-independent stage in the stroma}{Calvin cycle}{D}
\qa{Water falling from clouds}{Precipitation}{E}
\qa{Gas added by burning fossil fuels}{\ce{CO2}}{E}
\begin{confbox}
\confusion{Photosynthesis}{respiration}{photosynthesis stores energy (uses \ce{CO2}, makes \ce{O2}); respiration releases it (uses \ce{O2}, makes \ce{CO2}). Plants do \emph{both}.}
\confusion{Aerobic}{anaerobic}{aerobic uses oxygen and yields much ATP; anaerobic has no oxygen and yields little ATP.}
\confusion{Evaporation}{transpiration}{transpiration is evaporation of water specifically from plant leaves.}
\end{confbox}

\topic{Development of the Atomic Model}
\kf
\begin{keyfacts}
\item \textbf{Democritus}: ``atomos'' (indivisible).
\item \textbf{Dalton} (early 1800s): solid-sphere (billiard-ball) model and atomic theory.
\item \textbf{J.J. Thomson} (1897): discovered the electron with a cathode-ray tube; plum-pudding model.
\item \textbf{Rutherford} (1911): gold-foil experiment showed a small, dense, positive \textbf{nucleus} and mostly empty space.
\item \textbf{Bohr} (1913): electrons in fixed energy levels (shells); planetary model.
\item Modern view: electron-cloud (quantum mechanical) model, \textbf{Schr\"odinger}.
\item \textbf{Chadwick} (1932) discovered the neutron.
\item Shells for the first 20 elements fill 2, 8, 8, 2.
\end{keyfacts}
\cards
\qa{Coined ``atomos''}{Democritus}{A}
\qa{Solid-sphere (billiard-ball) model}{John Dalton}{E}
\qa{Plum-pudding model}{J.J. Thomson}{E}
\qa{Particle found with the cathode-ray tube}{Electron}{A}
\qa{Gold-foil experiment}{Ernest Rutherford}{E}
\qa{Main conclusion of the gold-foil experiment}{Dense positive nucleus; mostly empty space}{A}
\qa{Electrons in fixed energy levels}{Niels Bohr}{E}
\qa{Discovered the neutron}{James Chadwick}{A}
\qa{Current model of the atom}{Electron-cloud (quantum) model}{A}
\qa{Maximum electrons in the 2nd shell}{8}{A}
\qa{Electron arrangement of sodium ($Z = 11$)}{2, 8, 1}{D}
\begin{confbox}
\confusion{Thomson}{Rutherford model}{Thomson: positive ``pudding'' with electrons, no nucleus. Rutherford: a tiny nucleus with electrons around it.}
\confusion{Bohr}{electron-cloud model}{Bohr places electrons in exact orbits; the cloud model gives only regions of probability.}
\end{confbox}

\topic{Subatomic Particles}
\kf
\begin{keyfacts}
\item \textbf{Proton} (\ce{p+}): charge $+1$, mass about 1\,amu, in the nucleus.
\item \textbf{Neutron} (\ce{n^0}): no charge, mass about 1\,amu, in the nucleus.
\item \textbf{Electron} (\ce{e-}): charge $-1$, mass about 1/1836 of a proton, in the shells.
\item Atomic number $Z$ = number of protons. Mass number $A$ = protons + neutrons. In a neutral atom, electrons = protons.
\end{keyfacts}
\cards
\qa{Charge of a neutron}{Zero (neutral)}{E}
\qa{Particle with the smallest mass}{Electron}{E}
\qa{Atomic number equals the number of}{Protons}{E}
\qa{Mass number equals}{Protons + neutrons}{E}
\qa{Where protons and neutrons are found}{Nucleus}{E}
\qa{About how many times heavier a proton is than an electron}{About 1836}{D}
\begin{confbox}
\confusion{Atomic number}{mass number}{atomic number counts protons only; mass number counts protons + neutrons.}
\confusion{Mass number}{atomic mass}{mass number is a whole count for one isotope; atomic mass is the decimal average shown on the table.}
\end{confbox}
\mnemonic{\textbf{P-E-N}: \textbf{P}roton \textbf{P}ositive, \textbf{E}lectron n\textbf{E}gative, \textbf{N}eutron \textbf{N}eutral.}

% ---------------------------------------------------------------
\subsection{Term 2: Elements, Periodic Table, Earth's Crust and Volcanoes}

\topic{Elements and Compounds}
\kf
\begin{keyfacts}
\item Pure substances have a fixed composition: \textbf{elements} and \textbf{compounds}.
\item Every atom of an element has the same, unique number of protons.
\item Compounds are two or more elements chemically combined in a fixed ratio, separable only by chemical means (e.g., electrolysis).
\item Mixtures have a variable composition and are homogeneous (solutions) or heterogeneous.
\end{keyfacts}
\cards
\qa{Substance with only one kind of atom}{Element}{E}
\qa{Water: element, compound, or mixture?}{Compound}{E}
\qa{What makes each element unique}{Its number of protons}{A}
\qa{Is air a pure substance?}{No, it is a mixture}{A}
\qa{Process that splits water into \ce{H2} and \ce{O2}}{Electrolysis}{A}
\qa{Ratio of H to O atoms in water}{2:1}{E}
\qa{Seawater is a}{Homogeneous mixture}{A}
\qa{Can a compound be separated by filtration?}{No, only by chemical means}{D}
\begin{confbox}
\confusion{Compound}{mixture}{a compound has a fixed ratio and new properties; a mixture has a variable ratio and its parts keep their properties.}
\confusion{Atom}{molecule}{a molecule is two or more atoms bonded together (\ce{O2}, \ce{H2O}).}
\end{confbox}

\topic{The Periodic Table}
\kf
\begin{keyfacts}
\item \textbf{D\"obereiner}: triads. \textbf{Newlands}: law of octaves. \textbf{Mendeleev} (1869): ordered by atomic mass, left gaps, and predicted gallium and germanium. \textbf{Moseley}: ordered by atomic number (the modern periodic law).
\item The table has 7 periods (rows) and 18 groups (columns).
\item Elements in the same group have the same number of valence electrons and similar properties. The period number is the number of occupied shells.
\item Group 1 alkali metals, Group 2 alkaline earth metals, Groups 3--12 transition metals, Group 17 halogens, Group 18 noble gases.
\item Latin-based symbols: Na (natrium), K (kalium), Fe (ferrum), Cu (cuprum), Ag (argentum), Au (aurum), Pb (plumbum), Hg (hydrargyrum), Sn (stannum).
\item Example: \ce{^{27}_{13}Al} has 13 protons, 13 electrons, and $27 - 13 = 14$ neutrons.
\end{keyfacts}
\cards
\qa{Law of triads}{Johann D\"obereiner}{A}
\qa{Law of octaves}{John Newlands}{A}
\qa{Father of the periodic table}{Dmitri Mendeleev}{E}
\qa{Arranged elements by atomic number}{Henry Moseley}{A}
\qa{Mendeleev's ``eka-aluminium'' is now}{Gallium}{D}
\qa{Number of periods}{7}{E}
\qa{Number of groups}{18}{E}
\qa{Group 18 elements}{Noble gases}{E}
\qa{Group 17 elements}{Halogens}{E}
\qa{Group 1 elements}{Alkali metals}{E}
\qa{Symbol of potassium}{K}{E}
\qa{Symbol of sodium}{Na}{E}
\qa{Element with symbol Au}{Gold}{E}
\qa{Element with symbol Hg}{Mercury}{E}
\qa{Latin name behind Fe}{Ferrum}{A}
\qa{Neutrons in aluminium-27}{14}{A}
\qa{Valence electrons of Group 2 elements}{2}{A}
\qa{Element in period 3, group 17}{Chlorine}{A}
\qa{Element number 20}{Calcium}{A}
\qa{Why noble gases are unreactive}{Full outer shell}{A}
\qa{Metalloid used in computer chips}{Silicon}{A}
\begin{confbox}
\confusion{Group}{period}{groups are columns (same valence electrons); periods are rows (same number of shells).}
\confusion{Mendeleev}{Moseley}{Mendeleev ordered by atomic \emph{mass}; Moseley ordered by atomic \emph{number}.}
\confusion{Metal}{metalloid}{metalloids (B, Si, Ge, As, Sb, Te) lie on the ``staircase'' and have mixed properties.}
\end{confbox}

\topic{Earth's Surface and Crust}
\kf
\begin{keyfacts}
\item About 71\% of Earth's surface is water and about 29\% is land.
\item \textbf{Oceanic crust} is thin, dense, basaltic, and young. \textbf{Continental crust} is thick, less dense, granitic, and old.
\item The upper crustal layers are \textbf{sial} (silica + alumina, continental) over \textbf{sima} (silica + magnesium, oceanic).
\item \textbf{Lithosphere} = crust + uppermost mantle. It is rigid and broken into plates.
\item Volcanoes, earthquake epicenters, and young mountain belts cluster along plate boundaries, such as the Pacific Ring of Fire.
\end{keyfacts}
\cards
\qa{Percent of Earth's surface covered by water}{About 71\%}{E}
\qa{Thinner, denser type of crust}{Oceanic crust}{A}
\qa{Main rock of oceanic crust}{Basalt}{A}
\qa{Main rock of continental crust}{Granite}{A}
\qa{Crust plus uppermost mantle}{Lithosphere}{A}
\qa{Silica--alumina-rich continental layer}{Sial}{D}
\qa{Where volcanoes, quakes, and mountains cluster}{Plate boundaries}{A}
\qa{Volcano and quake belt around the Pacific}{Ring of Fire}{E}

\topic{Volcanoes}
\kf
\begin{keyfacts}
\item By activity (PHIVOLCS): \textbf{active}, \textbf{potentially active}, \textbf{inactive}. The PHIVOLCS active list has \verify{24} volcanoes.
\item \textbf{Shield}: broad, gentle slopes, fluid basaltic lava. \textbf{Cinder cone}: small and steep, made of pyroclastics. \textbf{Composite (strato)}: lava and ash layers, steep and symmetrical (Mayon).
\item Silica content controls viscosity: high silica gives thick magma and explosive eruptions; low silica gives runny lava and gentle eruptions.
\item Eruption types: Hawaiian, Strombolian, Vulcanian, Plinian (most violent; Pinatubo 1991), and phreatic (steam-driven; common at Taal).
\item A \textbf{caldera} is a large basin formed when a volcano collapses (Taal Lake).
\item PH examples: Mayon (Albay), Taal (Batangas), Pinatubo (Zambales), Kanlaon (Negros), Bulusan (Sorsogon), Hibok-Hibok (Camiguin).
\end{keyfacts}
\cards
\qa{Volcano type of Mayon}{Composite (stratovolcano)}{E}
\qa{Broad, gently sloping volcano of fluid lava}{Shield volcano}{E}
\qa{Small, steep cone of cinders}{Cinder cone}{E}
\qa{Large basin from a collapsed volcano}{Caldera}{A}
\qa{Steam-driven eruption}{Phreatic}{A}
\qa{Most violent eruption type, with a tall column}{Plinian}{A}
\qa{Year Pinatubo erupted}{1991}{E}
\qa{Magma property that makes eruptions explosive}{High viscosity (high silica)}{D}
\qa{Low-silica lava is}{Fluid (runny)}{A}
\qa{Volcano in a lake in Batangas}{Taal}{E}
\qa{Active volcano on Negros}{Kanlaon}{A}
\qa{Active volcano on Camiguin}{Hibok-Hibok}{D}
\qa{Number of active PH volcanoes (PHIVOLCS)}{\verify{24}}{D}
\begin{confbox}
\confusion{Magma}{lava}{magma is molten rock underground; lava is molten rock at the surface.}
\confusion{Crater}{caldera}{a crater is a small vent opening; a caldera is a large collapse basin.}
\confusion{Oceanic}{continental crust}{oceanic is thin, dense, and basaltic; continental is thick, light, and granitic.}
\end{confbox}

% ---------------------------------------------------------------
\subsection{Term 3: Typhoons, Tides, Acceleration, Energy, Light}

\topic{Typhoons}
\kf
\begin{keyfacts}
\item A tropical cyclone is a low-pressure system that forms over warm ocean water (about \SI{26.5}{\celsius} or warmer). It rotates counterclockwise in the Northern Hemisphere (Coriolis effect).
\item PAGASA (under DOST) categories, by sustained winds: TD $\le$ \SI{61}{km/h}; TS 62--88; STS 89--117; TY 118--184; Super Typhoon $\ge$ \SI{185}{km/h} (since 2022).
\item Warning signals: TCWS No.\ 1 to No.\ 5.
\item Parts: the \textbf{eye} is calm with the lowest pressure; the \textbf{eyewall} has the strongest winds; then the rainbands.
\item A typhoon weakens over land and mountains because it loses its warm-water fuel. The Sierra Madre shields eastern Luzon.
\item About 20 tropical cyclones enter the \textbf{PAR} each year. PAR corners: 5\textdegree N 115\textdegree E, 15\textdegree N 115\textdegree E, 21\textdegree N 120\textdegree E, 25\textdegree N 120\textdegree E, 25\textdegree N 135\textdegree E, 5\textdegree N 135\textdegree E.
\item A \textbf{storm surge} is an abnormal rise of the sea (Yolanda, 2013, Tacloban).
\end{keyfacts}
\cards
\qa{PH agency issuing typhoon bulletins}{PAGASA}{E}
\qa{PAGASA is under which department}{DOST}{A}
\qa{Calmest part of a typhoon}{Eye}{E}
\qa{Part with the strongest winds}{Eyewall}{A}
\qa{Spin of a typhoon in the Northern Hemisphere}{Counterclockwise}{A}
\qa{Effect that deflects winds and makes storms spin}{Coriolis effect}{D}
\qa{Why typhoons weaken after landfall}{Cut off from warm water}{A}
\qa{Mountain range shielding eastern Luzon}{Sierra Madre}{A}
\qa{Abnormal sea rise pushed by typhoon winds}{Storm surge}{E}
\qa{Year of Super Typhoon Yolanda}{2013}{A}
\qa{Highest TCWS number}{Signal No.\ 5}{A}
\qa{Typhoons entering the PAR per year (average)}{About 20}{A}
\qa{Minimum wind for a PAGASA super typhoon}{\SI{185}{km/h}}{D}
\qa{PAGASA wind range for ``typhoon''}{118--\SI{184}{km/h}}{D}
\qa{Eastern boundary of the PAR}{135\textdegree E}{D}
\qa{Approximate minimum sea temperature for formation}{About \SI{26.5}{\celsius}}{D}
\begin{confbox}
\confusion{Eye}{eyewall}{the eye is calm; the eyewall around it is the most violent part.}
\confusion{Signal number}{category}{the TCWS signal is the wind threat to \emph{your area}; the category is the storm's own strength.}
\end{confbox}

\topic{Tides}
\kf
\begin{keyfacts}
\item Tides are the regular rise and fall of sea level, caused mainly by the Moon's gravity and partly by the Sun's (about half as strong).
\item \textbf{Spring tides} (largest range) happen when the Sun, Earth, and Moon line up, at new and full moon.
\item \textbf{Neap tides} (smallest range) happen when the Sun and Moon are at right angles, at the quarter moons.
\item Most coasts get 2 high and 2 low tides a day.
\item Tidal energy is renewable. La Rance (France) is a classic tidal power plant.
\end{keyfacts}
\cards
\qa{Main cause of tides}{Moon's gravity}{E}
\qa{Highest highs and lowest lows}{Spring tides}{E}
\qa{Tide at the quarter moons}{Neap tide}{A}
\qa{Moon phases with spring tides}{New and full moon}{A}
\qa{High tides per day on most coasts}{Two}{A}
\qa{Sun's tidal effect compared with the Moon's}{About half}{D}
\qa{Famous tidal power station in France}{La Rance}{D}
\begin{confbox}
\confusion{Spring}{neap tide}{spring = aligned, big range; neap = right angles, small range. ``Spring'' has nothing to do with the season.}
\end{confbox}

\topic{Acceleration and Motion Graphs}
\kf
\begin{keyfacts}
\item $a = \dfrac{v_f - v_i}{t}$, in \si{\meter\per\second\squared}. Forces cause acceleration.
\item An object accelerates when its speed \emph{or} direction changes. Uniform circular motion is acceleration toward the centre (centripetal).
\item Distance--time graph: a curve (parabola) means constant acceleration.
\item Velocity--time graph: the slope is the acceleration; the area under it is the displacement; a flat line means constant velocity.
\item Free-fall acceleration $g \approx \SI{9.8}{\meter\per\second\squared}$.
\end{keyfacts}
\cards
\qa{SI unit of acceleration}{\si{\meter\per\second\squared}}{E}
\qa{0 to \SI{20}{\meter\per\second} in \SI{5}{s}}{\SI{4}{\meter\per\second\squared}}{A}
\qa{Slope of a velocity--time graph}{Acceleration}{A}
\qa{Area under a velocity--time graph}{Displacement}{D}
\qa{Horizontal line on a v--t graph}{Constant velocity (zero acceleration)}{A}
\qa{d--t graph shape for constant acceleration}{Curve (parabola)}{D}
\qa{Car rounding a curve at constant speed: accelerating?}{Yes, its direction changes}{A}
\qa{Direction of acceleration in uniform circular motion}{Toward the centre}{D}
\qa{Acceleration due to gravity}{\SI{9.8}{\meter\per\second\squared}}{E}
\qa{Slowing down is called}{Deceleration}{E}
\begin{confbox}
\confusion{Flat d--t line}{flat v--t line}{a flat d--t line means at rest; a flat v--t line means constant velocity.}
\confusion{Velocity}{acceleration}{at the top of a vertical throw $v = 0$ but $a = g$ is still acting.}
\end{confbox}

\topic{Work, Energy, Power, and Hydroelectricity}
\kf
\begin{keyfacts}
\item Kinetic energy $KE = \frac{1}{2}mv^2$. Gravitational potential energy $PE = mgh$. Both are in joules.
\item Work $W = Fd$ (force along the displacement). No displacement, or force $\perp$ motion, means no work.
\item Power $P = W/t$, in watts. 1\,hp $\approx$ \SI{746}{W}.
\item Mechanical energy = $KE + PE$. Energy is conserved: a falling object turns PE into KE.
\item Hydroelectric plants turn the PE of dam water into KE, which turns a turbine and generator to make electricity. Examples: Angat (Bulacan; supplies most of Metro Manila's water), San Roque, Magat, Pantabangan.
\end{keyfacts}
\cards
\qa{Energy of motion}{Kinetic energy}{E}
\qa{Stored energy due to height}{Gravitational potential energy}{E}
\qa{Unit of work and energy}{Joule}{E}
\qa{Formula for work}{$W = Fd$}{E}
\qa{\SI{10}{N} pushes a box \SI{3}{m}: work}{\SI{30}{J}}{A}
\qa{Pushing a wall that doesn't move: work done}{Zero}{A}
\qa{Rate of doing work}{Power}{E}
\qa{Unit of power}{Watt}{E}
\qa{KE + PE}{Mechanical energy}{A}
\qa{Energy that is maximum at the top of a swing}{Potential energy}{A}
\qa{Energy change in a falling coconut}{PE $\rightarrow$ KE}{E}
\qa{Energy is neither created nor destroyed}{Law of Conservation of Energy}{E}
\qa{Doubling speed multiplies KE by}{4}{D}
\qa{1 horsepower is about}{\SI{746}{W}}{D}
\qa{Carrying a bag while walking level at constant speed: work done on the bag}{Zero (force $\perp$ motion)}{D}
\qa{Bulacan dam supplying most of Metro Manila's water}{Angat Dam}{A}
\qa{Hydropower changes water's PE into}{Electrical energy (via KE)}{A}
\begin{confbox}
\confusion{Energy}{power}{energy is the total amount (J); power is how fast it is used or transferred (W).}
\confusion{Work (physics)}{effort}{holding a heavy box still is tiring, but the work done on it is zero.}
\end{confbox}

\topic{Properties of Light: Reflection and Refraction}
\kf
\begin{keyfacts}
\item Law of reflection: angle of incidence = angle of reflection, both measured from the normal.
\item Plane mirror image: virtual, upright, same size, laterally inverted.
\item \textbf{Concave} (converging) mirrors are used in flashlights, solar cookers, and shaving mirrors. \textbf{Convex} (diverging) mirrors give a wide view and are used as vehicle side mirrors and store mirrors.
\item \textbf{Refraction} is light bending as its speed changes between media. Entering a denser medium (air $\rightarrow$ water), it bends toward the normal.
\item A prism disperses white light into ROYGBIV; violet bends the most.
\item A convex lens is a magnifier and corrects farsightedness. A concave lens corrects nearsightedness.
\item Total internal reflection makes optical fibres work.
\end{keyfacts}
\cards
\qa{Law of reflection}{Angle of incidence = angle of reflection}{E}
\qa{Plane mirror image: real or virtual?}{Virtual}{E}
\qa{``AMBULANCE'' is written in reverse because mirrors cause}{Lateral inversion}{A}
\qa{Mirror used for vehicle side mirrors}{Convex}{E}
\qa{Mirror in flashlights and solar cookers}{Concave}{A}
\qa{Bending of light between media}{Refraction}{E}
\qa{Splitting white light into colours}{Dispersion}{A}
\qa{Colour bent most by a prism}{Violet}{D}
\qa{Lens of a magnifying glass}{Convex}{E}
\qa{Lens that corrects nearsightedness}{Concave}{D}
\qa{Light going from air into water bends}{Toward the normal}{D}
\qa{Principle behind optical fibres}{Total internal reflection}{D}
\qa{Speed of light in a vacuum}{\SI{3e8}{\meter\per\second}}{A}
\begin{confbox}
\confusion{Reflection}{refraction}{reflection bounces light back; refraction bends it as it passes through.}
\confusion{Concave}{convex mirror}{concave ``caves in'' and converges light; convex bulges out, diverges, and widens the view.}
\confusion{Real}{virtual image}{a real image can appear on a screen (light actually meets); a virtual image cannot.}
\end{confbox}
\mnemonic{\textbf{ROY G BIV}: Red, Orange, Yellow, Green, Blue, Indigo, Violet (longest to shortest wavelength).}

\end{gradebody}
\cardstats
